\documentclass[../DoAn.tex]{subfiles}
\begin{document}

\begin{center}
    \Large{\textbf{TÓM TẮT NỘI DUNG ĐỒ ÁN}}\\
\end{center}
\vspace{1cm}

Trong bối cảnh nhu cầu học tiếng Nhật ngày càng tăng, người học thường phải sử dụng nhiều công cụ rời rạc như từ điển, ứng dụng flashcard, tài liệu ngữ pháp, công cụ dịch và nền tảng luyện giao tiếp. Cách học phân tán này khiến người học khó theo dõi tiến độ, khó ôn tập đúng thời điểm và chưa tận dụng hiệu quả dữ liệu lỗi sai cá nhân. Các công cụ hiện có đã hỗ trợ tốt từng nhu cầu riêng lẻ, nhưng còn hạn chế trong việc kết nối dữ liệu học tập, cá nhân hóa nội dung ôn tập và giải thích nghĩa theo ngữ cảnh chuyên ngành.

Để giải quyết vấn đề trên, đồ án lựa chọn hướng tiếp cận xây dựng một hệ thống học tiếng Nhật thông minh, tích hợp các hoạt động tra cứu, ghi nhớ, ôn tập, luyện hội thoại và dịch thuật trong cùng một nền tảng. Hướng tiếp cận này được lựa chọn vì có thể tạo ra vòng học tập liên tục, trong đó dữ liệu từ quá trình tra cứu, flashcard, lỗi sai và phiên luyện tập được sử dụng lại để hỗ trợ người học hiệu quả hơn.

Giải pháp của đồ án gồm các chức năng chính: tra cứu từ vựng, kanji và ngữ pháp; quản lý bộ thẻ flashcard; ôn tập theo thuật toán lặp lại ngắt quãng SM-2; sinh trắc nghiệm và câu chuyện tương tác dựa trên bộ thẻ, cấp độ JLPT và lỗi sai gần đây; luyện hội thoại với AI Tutor qua văn bản hoặc giọng nói; và dịch chuyên ngành tiếng Nhật với sự hỗ trợ của đồ thị tri thức Neo4j kết hợp mô hình ngôn ngữ lớn.

Đóng góp chính của đồ án là thiết kế và triển khai một nguyên mẫu hệ thống học tiếng Nhật có khả năng kết nối frontend, backend, cơ sở dữ liệu và tầng AI trong một sản phẩm thống nhất. Kết quả đạt được là một hệ thống có thể hỗ trợ người học tra cứu, ghi nhớ, ôn tập, luyện giao tiếp và nhận phản hồi thông minh, tạo nền tảng để mở rộng thành trợ lý học ngoại ngữ cá nhân hóa trong tương lai.

\begin{flushright}
Sinh viên thực hiện\\
\begin{tabular}{@{}c@{}}
\textit{(Ký và ghi rõ họ tên)}
\end{tabular}
\end{flushright}

\end{document}
